刚体是忽略形变的理想质点系：任意两质点之间的距离在运动中保持不变。真实固体受力后都会
发生形变，只有形变相对物体尺度很小、且对所研究运动的影响可忽略时，才采用刚体模型。
它不是“绝对坚硬”的真实材料；局部压缩、弹性波、内部振动及断裂需要另行考虑。

\subsection{位置、姿态与自由度}
质心的三个坐标只确定位置，不确定朝向。例如一本书可以质心不动而转动。自由三维刚体需要
三个位置参数和三个姿态参数，共六个构型自由度。固定一点后剩三个转动自由度；固定一条轴后
剩一个；平面运动有两个平移和一个转动自由度。无滑动接触还在速度层面增加约束，不能重复计数。

欧拉角（Euler angles）是一组描述姿态的坐标。取原点都在质心的两组右手轴：xyz的方向
保持相对惯性系不变，XYZ与刚体固连。前者若随质心加速平移，其原点并非惯性系中的固定点，
“方向固定”不等于“整个坐标系是惯性系”。采用 $z$--$N$--$Z$ 转动约定：
\begin{enumerate}[label=(\arabic*),leftmargin=2.4em]
 \item 绕z转$\alpha$，把x转到节线N；N是xy与XY平面的交线，取向由转动约定给出。
 \item 绕N按右手规则转$\beta$，把z转到Z。
 \item 绕Z转$\gamma$，把N转到固连轴X。
\end{enumerate}
图~\ref{ch09:fig:euler}把三次转动分开画；第二面板N指向纸内，因而正$\beta$在该观察方向下是顺时针。
三面板的视向不同，不能把它们的纸面左右直接当成同一组三维坐标。
\eulerDiagram
通常约定 $0\le\alpha,\gamma<2\pi$、$0\le\beta\le\pi$。
当$\beta$为0或$\pi$时，两平面重合，节线不唯一，$\alpha$与$\gamma$不能独立确定；这是坐标退化，不是刚体失去姿态。
因此欧拉角只在排除退化并固定角度范围后局部唯一。位置和姿态确定\textbf{构型}；要确定
动力学状态，还须给质心速度和角速度，或等价的广义速度。

有限转动一般不可交换，先绕一轴再绕另一轴与倒序结果不同。无穷小转动在一阶近似下才可按
矢量相加。故实际角速度应为
\begin{equation}
 \vect\omega=\dot\alpha\,\uvec e_z+\dot\beta\,\uvec e_N+\dot\gamma\,\uvec e_Z,
\end{equation}
右侧三根轴不通常互相垂直，应投影到同一组轴再求分量；三个角的导数不是三个固定直角分量。
定轴或平面运动只需一个角$\varphi$时，才可直接以有符号量写 $\omega=\dot\varphi$、$\beta=\ddot\varphi$。

\subsection{刚体速度、加速度与瞬心}
平动中$\omega$=0，各点速度相同，任一固连线段始终平行于自身；定轴转动中固连轴不动，
各点绕该轴作圆周运动；平面平行运动则包含质心平动及绕质心的转动，车轮和滑杆是常见模型。
对刚体任意两点A、B，有
\begin{equation}
 \vect v_B=\vect v_A+\vect\omega\times\overrightarrow{AB}.\label{ch09:eq:rigid-velocity}
\end{equation}
与 $\overrightarrow{AB}$ 点乘，可得
$(\vect v_B-\vect v_A)\cdot\overrightarrow{AB}=0$，即相对速度在连线方向上的分量为零。
这也等于对固定长度条件 $\overrightarrow{AB}\cdot\overrightarrow{AB}=\text{常量}$ 求导。

平面运动且$\omega$不为零时，速度场可等效为绕瞬时速度为零的点P转动。若已知两点速度方向，
过两点作各自速度的垂线，其非平行交点给出瞬心；平行等速的特殊情形应回到速度方程判断，
不能强求一个有限交点。瞬心可以位于刚体之外的延伸速度场中，不一定是实际材料点。
对非零速度纯平动，$\omega$=0，不存在有限距离的瞬时静止点。

图~\ref{ch09:fig:rolling}中车轮向右纯滚，接触地面的P瞬时静止，故任一点Q满足
$\vect v_Q=\vect\omega\times\overrightarrow{PQ}$。
速度必须垂直PQ，且 $v_Q/\omega=PQ$。顶点 $P'$ 到P的距离为2R，故顶点速率为质心速率的两倍。
\rollingKinematicsDiagram
瞬心不是一个始终固定的铰点。下一时刻轮缘接触地面的是另一个材料点，所以不能把整个过程
当成绕同一P的定轴转动。对同一对材料点，式~\eqref{ch09:eq:rigid-velocity}求导给出
\begin{equation}
 \vect a_B=\vect a_A+\vect\beta\times\overrightarrow{AB}
 +\vect\omega\times(\vect\omega\times\overrightarrow{AB}).
\end{equation}
第二项是切向项，第三项是法向项。P瞬时速度为零，不代表其加速度为零。
例如水平面上以恒定质心速度纯滚的圆轮，接地点材料点有竖直向上的加速度 $\omega^2R$。
因此瞬心法可用于速度，不能未经推导就用 $\sum M_P=I_P\beta$ 代替质心动力学方程。

\subsection{刚体的动量与角动量}
将各质元速度相加，利用相对质心速度的质量加权和为零，得到
\begin{equation}
 \vect P=M\vect v_C,\qquad \sum\vect F_{\rm ext}=M\vect a_C.
\end{equation}
同一外力无论作用在何处，对质心加速度的贡献相同；但对转动的作用由其作用线决定。
角动量则为 $\vect L_O=\overrightarrow{OC}\times M\vect v_C+\vect L_C$。
前一节已说明 $\vect L_C$ 不一定与 $\vect\omega$ 平行，图~\ref{ch09:fig:nonparallel}给出具体反例。
\nonparallelAngularDiagram
设轻杆两端各质量m，质心O到各端距离l，杆与竖直轴夹角$\alpha$。图中 $+x$ 向右、$+z$ 向上，
为右手系，$+y$ 指向纸内。若 $\vect\omega=(0,0,\omega)$，上端A满足
\begin{equation}
 \vect r_A=(l\sin\alpha,0,l\cos\alpha),\qquad
 \vect v_A=(0,\omega l\sin\alpha,0).
\end{equation}
下端B的位置、速度均相反，所以两端对O的角动量相同，求和为
\begin{equation}
 \vect L_O=2m\vect r_A\times\vect v_A
 =2ml^2\omega(-\sin\alpha\cos\alpha,0,\sin^2\alpha).
\end{equation}
当 $0<\alpha<\pi/2$，它既有负x分量又有正z分量，不与竖直角速度同向。
其轴向分量仍是 $L_z=(2ml^2\sin^2\alpha)\omega=I_z\omega$。
这说明标量轴向关系可以正确，而向量式 $\vect L=I\vect\omega$ 却不成立。

一般关系为 $\vect L_C=\mathbf I_C\vect\omega$。仅沿主惯量轴转动时，角动量与角速度
平行，才能用一个标量惯量描述整个向量关系；球对称物体每个过中心方向都是主轴。
一般自由三维转动即使外力矩为零，L保持固定，刚体的轴仍可能摆动。

对平面刚体，最常用的方程组是
\begin{equation}
 \sum F_x=Ma_{Cx},\qquad \sum F_y=Ma_{Cy},\qquad \sum M_C=I_C\beta.
\end{equation}
最后一式只取垂直运动平面的分量，转轴、力矩和$\beta$使用同一正方向。若物体确绕惯性系中的固定轴
转动，也可用关于该轴的惯量写 $\sum M=I\beta$。对任意运动点则须使用前一节的参考点变换式，
不能因为某点瞬时静止就视为始终固定。摩擦、支持力通常是联立方程的未知量，不应先凭直觉定值。

\subsection{斜面纯滚动与摩擦判据}
设质量M、外半径R、绕质心轴惯量I的圆形刚体从静止沿倾角$\theta$的固定斜面下滚，
无外加驱动转矩。沿坡向下及顺时针为正，先假设静摩擦沿坡向上。
图~\ref{ch09:fig:incline}中P为接触点，$CP=R$ 且CP垂直斜面；N沿CP，重力通过C。
\inclineRollingDiagram
法向没有质心加速度，切向和转动方程分别为
\begin{equation}
 N=Mg\cos\theta,\qquad Mg\sin\theta-f=Ma_C,\qquad fR=I\beta.
\end{equation}
若不滑动，$a_C=\beta R$，因此 $f=Ia_C/R^2$，代回切向方程：
\begin{equation}
 Mg\sin\theta=(M+I/R^2)a_C,
 \qquad \boxed{a_C=\frac{g\sin\theta}{1+I/(MR^2)}}.
\end{equation}
继而求出实际所需的静摩擦
\begin{equation}
 f=\frac{I}{I+MR^2}Mg\sin\theta,\qquad
 \boxed{\mu_s\ge\frac fN=\frac{I}{I+MR^2}\tan\theta}.
\end{equation}
只有满足不等式，纯滚动假设才自洽。否则改用动摩擦 $\mu_kN$，并撤去 $a_C=\beta R$，
分别求平动与转动，直到接触点相对速度归零后再判断能否维持纯滚动。

令 $\kappa=I/(MR^2)$，图中实心球、实心圆柱、薄壁圆筒分别有
$\kappa=2/5,1/2,1$，相应加速度为 $5g\sin\theta/7,2g\sin\theta/3,g\sin\theta/2$。
在同一坡面、相同质心沿坡路程、同样从静止出发且均不滑时，实心球先到，薄壁圆筒后到。
总质量与绝对尺寸从比值中约去，决定先后的是质量分布。

静摩擦方向并非永远向上坡：加外转矩、改变作用力位置或让接触面运动，都可能使其反向或为零。
应先判断没有摩擦时接触点的相对滑动趋势，再用方程的正负校验。固定面上纯滚时
$P_f=\vect f\cdot\vect v_P=0$，但摩擦仍同时改变平动和转动动能；这两部分的功率相加为零。
运动传送带上的接触点对地速度不为零，静摩擦对物体可以做非零功。

\subsection{动能、外力功率与转动功}
任一质元速度分解为 $\vect v_i=\vect v_C+\vect\omega\times\vect r_i'$。
平方并求和，混合项含 $\sum_i m_i\vect r_i'=0$，因此
\begin{equation}
 E_k=\frac12Mv_C^2+\frac12\sum_i m_i|\vect\omega\times\vect r_i'|^2.
\end{equation}
平面刚体的角速度沿同一法向，第二项成为 $I_C\omega^2/2$；一般三维情形为
$\vect\omega\cdot\mathbf I_C\vect\omega/2$。不能把一般三维运动也一律套单个I。

外力作用点速度不同，应先逐力计算功率。利用标量三重积恒等式，
\begin{align}
 P&=\sum_i\vect F_i\cdot(\vect v_C+\vect\omega\times\vect r_i')\\
 &=\left(\sum_i\vect F_i\right)\cdot\vect v_C+
 \left(\sum_i\vect r_i'\times\vect F_i\right)\cdot\vect\omega\\
 &=\vect F_{\rm ext}\cdot\vect v_C+\vect M_C\cdot\vect\omega.
\end{align}
第一项给质心平动动能变化率，第二项给转动动能变化率。同一外力可同时参与两部分，
不能在计算功率时任意把作用点搬到质心。理想刚性约束的成对内力不对整体做净功，
但能在部分物体之间传递能量；若含主动内部驱动或可变形自由度，需另计相应功和储能。

绕固定轴转动，作用点的切向位移为 $r_\perp\dd\varphi$，故
\begin{equation}
 \dd W=M_{\rm axis}\dd\varphi,\qquad P=M_{\rm axis}\omega,
 \qquad W=\int_{\varphi_1}^{\varphi_2}M_{\rm axis}(\varphi)\dd\varphi=\Delta E_k.
\end{equation}
恒力矩且I不变才对应恒角加速度；力臂随姿态改变时应积分，不能简单用初始力矩乘总转角。
对纯滚动下坡，也可令势能减少等于两部分动能之和；它与动力学结果相符，且说明把物体当作
仅平动的质点会高估下坡速度和加速度。

\subsection{碰撞、打击与截面内力}
短时碰撞应同时考虑线冲量和冲量矩，对固定矩心O有
\begin{equation}
 M(\vect v_C^+-\vect v_C^-)=\vect J_{\rm ext},\qquad
 \vect L_O^+-\vect L_O^-=\int\vect M_O^{\rm ext}\dd t.
\end{equation}
若只有某一接触冲量J，对质心还有
$\Delta\vect L_C=\overrightarrow{CP}\times\vect J$；平面情况下可再写成I乘角速度改变量。
选择固定支点或碰撞点取矩，常能消去该处冲量矩，但不一定能消去其余外冲量矩。
打击中心就是平动与转动响应恰好配合，使固定支点不需提供冲量的打击位置；并非简单的几何中点。

碰撞通常伴随形变、声、热，动量或角动量守恒不保证动能守恒。仅有法向恢复系数等于1时，
还须查看切向摩擦是否耗能，不能直接宣称全部动能不变。
生鸡蛋内部液体引入相对流动自由度，不是单一角速度刚体；既要考虑流体参与转动的程度，
也要考虑耗散，后面的生熟鸡蛋题说明这一模型边界。

刚体假设各点距离不变，不意味着内部作用为零。杆或烟囱靠轴力、剪力和弯矩维持整体运动。
分析方法是先用整体方程求运动，再在待考察位置截开，对其中一段补齐截面相互作用，分别列力
与力矩方程。具体题里某项可以为零；只有实际非零的内力分布才用于比较断裂位置。
能量法适合连接状态、求角速度，动力学适合求瞬时反力及加速度，两者配合并以量纲、方向和
极限情况相互校验，不能仅靠整体能量判断哪个截面最危险。


\Needspace{0.34\textheight}
\begin{origexamplebox}{例题 15.1【砂轮制动】}
一个质量为 $m$、半径为 $R$ 的匀质砂轮以角速度 $\omega_0$ 旋转。在距离圆心为 $r$ 处
有恒力 $f$ 作用，方向如图。求随后砂轮的运动。
\brakeWheelDiagram
\end{origexamplebox}

\begin{xesolution}
图~\ref{ch09:fig:brake}中砂轮初始顺时针转动，f在圆心右侧向上，故力矩逆时针。
取顺时针为正，匀质砂轮视为均匀圆盘，绕中心轴的转动惯量
\[
 I=\frac12mR^2.
\]
力对轴的力矩大小为 $fr$，方向与初始转动相反，所以角加速度为常量
\[
 \beta=-\frac{fr}{I}=-\frac{2fr}{mR^2}.
\]
于是制动阶段
\[
 \boxed{\omega(t)=\omega_0-\frac{2fr}{mR^2}t},
 \qquad
 \boxed{\varphi(t)=\omega_0t-\frac{fr}{mR^2}t^2}.
\]
砂轮在
\[
 \boxed{t_s=\frac{mR^2\omega_0}{2fr}}
\]
时停止，停止前转过的角度为
\[
 \boxed{\varphi_s=\frac{mR^2\omega_0^2}{4fr}}.
\]
若题意中的恒力在停止后仍保持原方向，砂轮会反向转动；“制动”通常只讨论首次停转之前。
\end{xesolution}

\Needspace{0.37\textheight}
\begin{origexamplebox}{例题 15.2【两个砂轮】}
两个匀质砂轮质量分别为 $m_1,m_2$，半径分别为 $R_1,R_2$，初始分别以角速度
$\omega_{10},\omega_{20}$ 按图示方向旋转。$t=0$ 时两轮接触，此后接触处有恒定摩擦力
$f$（仅在相对滑动期间取此大小，两轮轴固定且忽略其他阻力矩）。求：
\begin{enumerate}[label=(\arabic*),leftmargin=2.2em]
  \item 任意时刻二者的角速度；
  \item 这个问题是否有类似的质点动力学问题相对应？
\end{enumerate}
\twoWheelsDiagram
\end{origexamplebox}

\begin{xesolution}
图~\ref{ch09:fig:two-wheels}的第一轮以顺时针为正、第二轮以逆时针为正；两接触表面初速度均向下。
令 $u_{10}=R_1\omega_{10}$、$u_{20}=R_2\omega_{20}$，以向下为正的初始相对速度
$s_0=u_{10}-u_{20}$。若 $s_0>0$，摩擦使第一轮减速、第二轮加速，图中力箭头对应此分支；
$s_0<0$ 时两摩擦箭头都反向。以 $\varepsilon=\operatorname{sgn}(s_0)$ 统一记号，
$I_i=\tfrac12m_iR_i^2$，持续滑动期间
\[
 \boxed{\omega_1(t)=\omega_{10}-\varepsilon\frac{fR_1}{I_1}t
 =\omega_{10}-\varepsilon\frac{2f}{m_1R_1}t},
\]
\[
 \boxed{\omega_2(t)=\omega_{20}+\varepsilon\frac{fR_2}{I_2}t
 =\omega_{20}+\varepsilon\frac{2f}{m_2R_2}t}.
\]
两接触表面不再相对滑动的条件为 $R_1\omega_1=R_2\omega_2$；因此滑动终止时刻
\[
 \boxed{t_*=
 \frac{|R_1\omega_{10}-R_2\omega_{20}|}
 {f(R_1^2/I_1+R_2^2/I_2)}}.
\]
上述含t的式子只用于 $0\le t\le t_*$。此后两轮无其他阻力矩，维持无相对滑动所需的静摩擦为零，
各自以最终角速度匀速转动。若初始 $s_0=0$，则 $t_*=0$，直接处于这一状态，不能强加大小f的动摩擦。

它对应于两个质点在恒定相互摩擦力下相对滑动的问题。映射关系是：切向线速度
$u_i=R_i\omega_i$ 对应质点速度，而“等效质量”
\[
 m_i^{\rm eff}=\frac{I_i}{R_i^2}
\]
对应质点质量。记 $q_i=I_i/R_i^2$，两轮切向方程为 $q_1\dot u_1=-\varepsilon f$、
$q_2\dot u_2=\varepsilon f$，相加得 $q_1u_1+q_2u_2$ 不变。于是最终共同轮缘速度及角速度为
\[
 u_* =\frac{q_1u_{10}+q_2u_{20}}{q_1+q_2},\qquad
 \omega_{1*}=\frac{u_*}{R_1},\quad \omega_{2*}=\frac{u_*}{R_2}.
\]
这不是说固定轴承系统的总线动量守恒，而是切向方程与两个相互作用质点的方程具有相同形式。
共同速度位于两初始轮缘速度之间；机械能损失
\[
 \Delta E=\frac12\frac{q_1q_2}{q_1+q_2}(u_{10}-u_{20})^2\ge0
\]
可由初末转动动能相减得到，量纲为能量且 $s_0=0$ 时为零。
\end{xesolution}

\Needspace{0.34\textheight}
\begin{origexamplebox}{例题 15.3【滚动的台球】}
质量为 $m$、半径为 $R$ 的匀质台球刚被击出时速度为 $v_0$，没有旋转。已知台球与地面的
滑动摩擦因数为 $\mu$，求台球此后的运动。
\billiardRollingDiagram
\end{origexamplebox}

\begin{xesolution}
图~\ref{ch09:fig:billiard}分别画出初态和滑动阶段，避免把 $\omega=0$ 的初态与后续转动混为一图。
初始接触点相对地面向前滑，滑动摩擦力大小为 $f=\mu mg$、方向向后。匀质实心球
$I=\tfrac25mR^2$。在滑动阶段
\[
 ma=-\mu mg,
 \qquad I\beta=\mu mgR.
\]
取使球进入向前纯滚动的角速度为正，得到
\[
 v(t)=v_0-\mu gt,
 \qquad
 \omega(t)=\frac{\mu mgR}{I}t=\frac{5\mu g}{2R}t.
\]
当 $v=\omega R$ 时开始无滑动滚动：
\[
 v_0-\mu gt_* =\frac52\mu gt_*,
 \qquad
 \boxed{t_* =\frac{2v_0}{7\mu g}}.
\]
相应速度和角速度为
\[
 \boxed{v_* =\frac57v_0},
 \qquad
 \boxed{\omega_* =\frac{5v_0}{7R}}.
\]
若忽略滚动阻力，$t>t_*$ 后台球以这两个恒定值作无滑动滚动。滑动阶段质心前进距离为
\[
 \boxed{x_* =v_0t_* -\frac12\mu gt_*^2
 =\frac{12v_0^2}{49\mu g}}.
\]
最终总动能为 $\tfrac12m(5v_0/7)^2+\tfrac12(2mR^2/5)[5v_0/(7R)]^2=5mv_0^2/14$，
比初值 $mv_0^2/2$ 少 $mv_0^2/7$。损失来自接触点滑动，不能用质心路程直接乘摩擦力代表总损失。
\end{xesolution}

\Needspace{0.35\textheight}
\begin{origexamplebox}{例题 15.4【台球碰撞】}
质量均为 $m$、半径均为 $R$ 的匀质白球和红球发生无摩擦正碰。碰前白球以速度 $v_0$
向前运动，同时具有角速度 $\omega_0$，且 $v_0=\omega_0R$。红球起初静止且不自转。
两球所在的同一水平面与球的动摩擦因数均为$\mu$，忽略滚动阻力；球间碰撞极短，可忽略该期间的地面冲量。
球间碰撞完全弹性。求碰撞后两个
球的运动，并判断白球是否可能再次撞到红球。
\ballsCollisionDiagram
\end{origexamplebox}

\begin{xesolution}
图~\ref{ch09:fig:balls}标出碰前状态，取向右与顺时针为正，碰后另以t=0计时。
球与球之间无摩擦，所以碰撞冲量沿两球心连线，不对任一球心产生力矩。因而碰撞过程中
两球角速度不变：
\[
 \omega_{\mathrm W}'=\omega_0,
 \qquad \omega_{\mathrm R}'=0.
\]
沿球心连线方向，这是相同质量的一维完全弹性碰撞，二球交换平动速度：
\[
 v_{\mathrm W}'=0,
 \qquad v_{\mathrm R}'=v_0.
\]
碰后白球仍自转，接触地面的点相对地面向后滑，所以地面对白球的滑动摩擦力向前；红球只
平动不转动，摩擦力向后。分别应用例 15.3 的滑动阶段公式：
\[
 v_{\mathrm W}(t)=\mu gt,
 \qquad
 \omega_{\mathrm W}(t)=\omega_0-\frac{5\mu g}{2R}t,
\]
\[
 v_{\mathrm R}(t)=v_0-\mu gt,
 \qquad
 \omega_{\mathrm R}(t)=\frac{5\mu g}{2R}t.
\]
两球分别达到纯滚动的时刻相同：
\[
 t_{\mathrm W}=t_{\mathrm R}=\frac{2v_0}{7\mu g}.
\]
且最终纯滚动速度为
\[
 \boxed{v_{\mathrm W}=\frac27v_0},
 \qquad
 \boxed{v_{\mathrm R}=\frac57v_0}.
\]
在整个滑动阶段 $v_{\mathrm R}-v_{\mathrm W}=v_0-2\mu gt>0$；两球进入纯滚动后红球仍更快，
所以间距单调增大，\boxed{\text{白球不会再次撞到红球}}。
更明确地，滑动阶段两球心间距减去2R的净间隙为
\[
 d(t)=\int_0^t(v_{\mathrm R}-v_{\mathrm W})\dd t
 =v_0t-\mu gt^2,\qquad 0\le t\le t_*.
\]
因此 $d(t_*)=10v_0^2/(49\mu g)>0$；之后间隙以 $3v_0/7$ 的恒定速率增大。
当$\mu$趋零时，滑动调节时间趋无穷，但球间碰后红球前进、白球只自转，仍不会追及。
\end{xesolution}

\Needspace{0.37\textheight}
\begin{origexamplebox}{例题 15.5【小明拍球】}
小明把一个篮球快速旋转地拍下。落地前篮球速度大小为 $v_0$、方向竖直向下，角速度为
$\omega_0$，转向如图。篮球视为质量 $m$、半径 $R$ 的空心球壳，与地面的滑动摩擦因数为
$\mu$，采用法向恢复系数为1、切向库仑摩擦的理想碰撞模型；一旦切向相对速度归零，
接触不再滑动，不考虑切向弹性回弹。求篮球弹起后的速度方向。
\bounceSpinProblemDiagram
\end{origexamplebox}

\begin{xesolution}
图~\ref{ch09:fig:bounce-analysis}将碰前状态、接触冲量与碰后方向分开。碰撞极短，重力冲量可忽略。
\bounceSpinDiagram
法向恢复系数为1意味着竖直速度反向而大小不变，因此
\[
 v_y'=v_0 \quad\text{（竖直向上）}.
\]
图中的逆时针自旋使落地点相对地面向右滑，地面对球的切向冲量 $J_t$ 向左。
先假设整个接触阶段都保持向右滑动，此时动摩擦冲量为
\[
 J_t=\mu J_n.
\]
法向冲量由竖直动量改变给出
\[
 J_n=m(v_0-(-v_0))=2mv_0,
\]
故水平速度
\[
 \boxed{v_x'=-\frac{J_t}{m}=-2\mu v_0}.
\]
于是弹起方向相对竖直线向左偏，偏角 $\alpha$ 满足
\[
 \boxed{\tan\alpha=\frac{|v_x'|}{v_y'}=2\mu}.
\]
同时空心球壳 $I=\tfrac23mR^2$，其自旋改变为
\[
 \omega' =\omega_0-\frac{J_tR}{I}
 =\omega_0-\frac{3\mu v_0}{R}.
\]
若碰撞全程保持滑动，即
\[
 \omega_0R>5\mu v_0.
\]
则以上结果成立。这个临界条件来自碰后接触点速度
$v_x'+\omega'R=\omega_0R-5\mu v_0$ 仍为正。
若 $\omega_0R\leq5\mu v_0$，摩擦足以使接触点在分离前或恰在分离时停止滑动。设所需切向
冲量为 $J_t$，由
\[
 v_x'+\omega'R
 =-\frac{J_t}{m}+\left(\omega_0-\frac{J_tR}{I}\right)R=0
\]
得
\[
 J_t=\frac{2}{5}m\omega_0R,
 \qquad
 v_x'=-\frac25\omega_0R.
\]
因此完整的弹起方向为
\[
 \boxed{
 \tan\alpha=
 \begin{cases}
  2\mu, & \omega_0R>5\mu v_0,\\[2pt]
  \dfrac{2\omega_0R}{5v_0}, & \omega_0R\leq5\mu v_0.
 \end{cases}}
\]
两种情形都向图中摩擦冲量所指的一侧偏转；第二种情形中 $\omega'=2\omega_0/5$，
与 $v_x'=-2\omega_0R/5$ 正好使接触点相对速度为零。
临界点两式相同；$\mu$=0时水平冲量为零，仍竖直弹起；$\omega_0=0$ 时也无侧向偏转。
法向动能恢复不等于总动能恢复。切向平动与转动的能量改变为
\[
 \Delta K_t=-J_t\omega_0R+\frac{J_t^2}{2}\left(\frac1m+\frac{R^2}{I}\right)\le0,
\]
在本题允许的冲量范围内非正，说明摩擦耗能与题设法向弹性条件相容。
\end{xesolution}

\Needspace{0.33\textheight}
\begin{origexamplebox}{例题 15.6【生熟鸡蛋】}
把不同鸡蛋放在同一个斜面上向下滚，如何利用物理方法区分生鸡蛋和熟鸡蛋？
\eggRollingDiagram
\end{origexamplebox}

\begin{xesolution}
图~\ref{ch09:fig:egg}把两只鸡蛋放在相同坡角的两个面板中，接触位置与释放高度相同；
用剖面区分内部凝固与内部液体。比较前还要尽量控制外形、质量、放置朝向和坡面的粗糙程度。
熟鸡蛋内部凝固，外壳与内部近似共同转动；生鸡蛋内部液体可以相对外壳流动，不能给全部质量
指定同一角速度。这个内部响应差异才是判断依据，而不是预先认定某一种一定先到。

先用圆形、质心位于几何中心的模型解释转动惯量的影响。如果外形在滚动截面近似为圆，
且凝固整体保持纯滚动，刚体近似给出
\[
 a_{\rm cooked}=\frac{g\sin\theta}{1+I/(mR^2)}.
\]
真实鸡蛋不是圆球，其接触半径和质心高度随姿态改变，上式只能作为圆形近似，不能精确描述
每个瞬间。生鸡蛋还同时存在两种不同影响：内部液体未充分随壳转动时，参与转动的有效惯量可能较小；
液体相对流动又通过黏滞作用耗散能量。两种影响不能只取其一。

例如取圆形壳内液体几乎不受黏性转矩的理想极限，平动质量为总质量m，参与转动的惯量主要是
外壳的 $I_{\rm shell}$，则可能有
\[
 a_{\rm liquid}\simeq\frac{g\sin\theta}{1+I_{\rm shell}/(mR^2)}.
\]
它反而可大于凝固整体的加速度。这不是对所有真实鸡蛋的预测，而是说明“存在黏滞耗能”
不能单独推出“生鸡蛋一定更慢”。在较强耦合极限，液体逐渐跟随外壳，又更接近共同转动；
真实结果还受形状和滚动姿态影响。

因此，斜面实验可记录滚动时间、姿态变化以及短暂改变外壳运动后的响应，并与已知生熟样品比较；
只给本题的条件，不能仅凭一次到达先后作无条件判别。为检验内部是否仍在运动，可另在水平桌面
将蛋旋转，短暂按停外壳后立即松手：生蛋内部液体若仍在转，会通过黏滞作用重新带动外壳；
熟蛋整体已停住，通常不会自行恢复转动。按住过久会让内部也停止，所以必须说明“短暂”。
这个补充实验利用内部相对运动，不能反过来当作斜面滚动谁先到底的证明。
\end{xesolution}

\Needspace{0.35\textheight}
\begin{origexamplebox}{例题 15.7【重滑轮】}
两个重物挂在滑轮两端，绳与定滑轮间不相对滑动。定滑轮质量为 $M$、半径为 $R$，质量均匀
分布。求松手时 $m_1$ 的加速度，设 $m_1>m_2$。
\massivePulleyDiagram
\end{origexamplebox}

\begin{xesolution}
图~\ref{ch09:fig:pulley}中绳绕轮缘且连到两重物的中心吊点，支架连接固定轮轴。
设 $m_1$ 向下、$m_2$ 向上加速度大小均为 $a$，两侧张力分别为 $T_1,T_2$：
\[
 m_1g-T_1=m_1a,
 \qquad
 T_2-m_2g=m_2a.
\]
均匀圆盘滑轮的转动惯量 $I=\tfrac12MR^2$。绳不打滑给出 $\beta=a/R$，滑轮转动方程为
\[
 (T_1-T_2)R=I\frac aR.
\]
联立消去两张力：
\[
 (m_1-m_2)g
 =\left(m_1+m_2+\frac I{R^2}\right)a.
\]
因此
\[
 \boxed{a=\frac{(m_1-m_2)g}{m_1+m_2+M/2}},
\]
方向为 $m_1$ 向下。滑轮有转动惯性，所以加速度小于把滑轮视为轻滑轮时的结果。
\end{xesolution}

\Needspace{0.39\textheight}
\begin{origexamplebox}{例题 15.8【力矩做功】}
两个质量均为 $m$ 的质点固定在长为 $2l$ 的轻杆上，一个在最远端，另一个在中点。杆的
另一端铰接在天花板上。初始时杆水平，由静止释放。求杆与水平呈 $\theta$ 角时：
\begin{enumerate}[label=(\arabic*),leftmargin=2.2em]
  \item 杆的角速度；
  \item 1 号、2 号质点的动能；
  \item 1 号、2 号质点重力势能的减小量；
  \item 说明本题与力矩做功的关系。
\end{enumerate}
\torqueWorkDiagram
\end{origexamplebox}

\begin{xesolution}
图~\ref{ch09:fig:torque-work}中$\theta$从水平参考线向下量，取第一次下落的 $0\le\theta\le\pi/2$。
两质点到转轴的距离分别为 $2l$ 与 $l$，系统转动惯量
\[
 I=m(2l)^2+ml^2=5ml^2.
\]
杆由水平转到向下倾斜 $\theta$ 时，两质点分别下降 $2l\sin\theta$ 与 $l\sin\theta$，
故机械能守恒给出
\[
 \frac12I\omega^2=mg(2l\sin\theta)+mg(l\sin\theta)=3mgl\sin\theta.
\]
所以
\[
 \boxed{\omega=\sqrt{\frac{6g\sin\theta}{5l}}}.
\]
两质点动能分别为
\[
 \boxed{E_{k1}=\frac12m(2l\omega)^2
 =\frac{12}{5}mgl\sin\theta},
\]
\[
 \boxed{E_{k2}=\frac12m(l\omega)^2
 =\frac35mgl\sin\theta}.
\]
两者重力势能减小量分别为
\[
 \boxed{-\Delta E_{p1}=2mgl\sin\theta},
 \qquad
 \boxed{-\Delta E_{p2}=mgl\sin\theta}.
\]
个别质点的动能增加并不等于它自己的重力势能减少，因为轻杆的内力在两质点之间传递能量；
但对整个刚体，支点不做功，总动能增加等于总重力势能减少。

从力矩做功看，重力对支点的合力矩为
\[
 M(\phi)=3mgl\cos\phi,
\]
故
\[
 W=\int_0^\theta M(\phi)\dd\phi
 =3mgl\sin\theta=\Delta E_k.
\]
这正是转动形式的动能定理。
\end{xesolution}

\Needspace{0.37\textheight}
\begin{origexamplebox}{例题 15.9【打击中心】}
一根质量均匀、质量为 $m$、长度为 $L$ 的杆悬挂在天花板上。用冲量 $J$ 水平打击杆上距
下端 $a$ 的位置，发现悬挂端没有受到冲量。求打击位置。
\percussionDiagram
\end{origexamplebox}

\begin{xesolution}
图~\ref{ch09:fig:percussion}中O是固定铰点，杆初始静止；打击点P距悬点为 $x=L-a$。
悬点冲量为零，杆质心获得的线动量完全来自打击冲量：
\[
 mv_C=J.
\]
对悬点取矩，悬点冲量没有力矩，角动量定理给出
\[
 I_O\omega=Jx,
 \qquad I_O=\frac13mL^2.
\]
刚体绕悬点转动，质心速度还满足
\[
 v_C=\omega\frac L2.
\]
代入前两式：
\[
 \frac Jm=\frac{Jx}{I_O}\frac L2,
\]
从而
\[
 \boxed{x=\frac{2L}{3}},
 \qquad
 \boxed{a=L-x=\frac L3}.
\]
即打击点距悬点 $2L/3$（距下端 $L/3$）。该点称为打击中心；在此处横向打击，悬点瞬间
不承受冲量。
\end{xesolution}

\Needspace{0.34\textheight}
\begin{origexamplebox}{例题 15.10【小明的疑惑】}
一根质量m、长度L、质量均匀分布的棍靠两端A、B的支撑点水平悬空放置。小明只分析左半根棍时，发现外力矩似乎
不能平衡。这是怎么回事？
\supportedRodProblemDiagram
\end{origexamplebox}

\begin{xesolution}
图~\ref{ch09:fig:rod-analysis}左图以整杆为研究对象，对称性与竖直力平衡给出两端支持力各为 $mg/2$，
合力与对任一点的合力矩均为零。
\supportedRodDiagram

但若只截取左半杆，除了左支点支持力 $mg/2$ 与左半杆重力 $mg/2$ 之外，截面还受到右半杆
施加的\textbf{内力和内力矩}。在本题中，中截面的水平、竖直力平衡分别给出轴力为零及
\[
 N_A-\frac{mg}{2}+Q=0\quad\Longrightarrow\quad Q=0.
\]
所以右图没有画一个实际非零的剪力箭头。两竖直外力大小相等、方向相反，但作用线不同，构成力偶。设杆长
为 $L$，左半杆重心距左支点 $L/4$，其外力偶矩大小为
\[
 M_{\rm ext}=\frac{mg}{2}\frac L4=\frac{mgL}{8}.
\]
以逆时针为正，左支持力与左半杆重力对中截面的合外力矩为
$-(mg/2)(L/2)+(mg/2)(L/4)=-mgL/8$。为保持静止，中截面必须提供逆时针弯矩，大小
\[
 \boxed{M_{\rm int}=\frac{mgL}{8}}
\]
的反向弯矩。

疑惑来自把刚体截开后只补了截面内力，却漏掉截面内力矩。对能承受弯曲的杆件，截面相互
作用一般包括轴力、剪力和弯矩；只有把它们全部画出，局部受力与力矩才能平衡。
\end{xesolution}

\Needspace{0.42\textheight}
\begin{origexamplebox}{例题 15.11【烟囱倒地】}
把烟囱视为质量 $m$、长度 $L$ 的均匀硬杆，绕底端倾倒，杆与竖直方向夹角为 $\theta$。
用两种模型判断哪里最容易断裂：
\begin{enumerate}[label=(\arabic*),leftmargin=2.2em]
  \item 烟囱内部会有垂直于烟囱的剪力，剪力最大处最先断裂；
  \item 烟囱内部会提供弯矩，弯矩最大处最先断裂。
\end{enumerate}
\chimneyDiagram
\end{origexamplebox}

\begin{xesolution}
图~\ref{ch09:fig:chimney}中 $x$ 从固定基底O沿杆向上量，研究截面P以上的一段；
规定切向正方向 $\uvec e_t$ 为倾倒的顺时针方向，角速度、角加速度也以顺时针为正。
线密度 $\lambda=m/L$。整杆绕底端转动惯量
$I_O=mL^2/3$。重力矩给出角加速度
\[
 \beta=\frac{mg(L/2)\sin\theta}{mL^2/3}
 =\frac{3g\sin\theta}{2L}.
\]
从竖直静止位置以任意小扰动启动的极限下，机械能守恒给出
\[
 \frac12I_O\omega^2=mg\frac L2(1-\cos\theta),
 \qquad
 \omega^2=\frac{3g}{L}(1-\cos\theta).
\]

在 $x$ 处截开，研究上段，长度 $s=L-x$、质量 $m_s=ms/L$，质心位于
$r_c=(L+x)/2$。沿垂直杆方向列质心运动方程。重力的垂杆分量为
$m_sg\sin\theta$，质心切向加速度为 $\beta r_c$。定义截面对上段的剪力为
$-Q\uvec e_t$，Q是有符号量，图中箭头表示Q的正向。因此
$m_sg\sin\theta-Q=m_s\beta r_c$，得
\begin{align*}
 Q(x)&=m_s\bigl(g\sin\theta-\beta r_c\bigr)\\
 &=\frac{mg\sin\theta}{4L^2}(L-x)(L-3x).
\end{align*}
Q为正时剪力与倾倒切向相反，为负时实际方向反转；剪力\textbf{大小}是 $|Q|$。令 $y=x/L$，$|Q|$ 正比于
$|(1-y)(1-3y)|$；在 $0\le y\le1$ 上最大值出现在 $y=0$，所以按模型一
\[
 \boxed{\text{基底附近最容易发生剪切断裂}}.
\]

再对上段质心取矩。重力通过上段质心，不产生力矩；截面剪力对上段质心产生力矩。为使上段
具有顺时针角加速度$\beta$，截面还必须提供弯矩。关于上段质心G，剪力 $-Q\uvec e_t$ 的
顺时针力矩是 $Qs/2$。定义截面弯矩的有符号量M以顺时针为正，则
\[
 M+Q\frac s2=I_c\beta,\qquad I_c=\frac{m_ss^2}{12}.
\]
移项并代入，可得
\begin{align*}
 M(x)&=I_c\beta-Q(x)\frac s2\\
 &=\frac{mg\sin\theta}{8L^2}(L-x)^3
   -\frac{mg\sin\theta}{8L^2}(L-x)^2(L-3x)\\
 &=\frac{mg\sin\theta}{4L^2}\,x(L-x)^2.
\end{align*}
\end{xesolution}

\begin{solutioncontinuation}
\enlargethispage{\baselineskip}
令 $y=x/L$，需最大化 $y(1-y)^2$：
\[
 \frac{\dd}{\dd y}\bigl[y(1-y)^2\bigr]
 =(1-y)(1-3y)=0.
\]
非端点极值为 $y=1/3$，所以按模型二
\[
 \boxed{x=\frac L3\quad\text{处弯矩最大，最容易弯折断裂}}.
\]
对 $0<\theta<\pi$、$0<x<L$，M为正，与图中的顺时针约定一致；两端M均为零。
Q具有力的量纲，M具有力乘长度的量纲，正负关系也满足上段转动方程。
值得注意的是，$M(x)$ 中的 $\sin\theta$ 只改变总体大小，不改变最大位置；以上结果针对题目
给定的理想均匀刚杆模型。
\end{solutioncontinuation}
